\chapter{Trabalhos relacionados}
\label{cap:trabalhos_relacionados}

O trabalho de \citeauthor{Dickens:08}~\cite{Dickens:08} introduziu o uso de MPI-IO
sobre o Lustre, identificando que as premissas de otimização do ROMIO conflitam com
o modelo de byte-range locking do filesystem: para padrões não-interleaved, nos quais
cada processo escreve em uma região exclusiva, escritas independentes evitam contenção
de locks e superam as coletivas em desempenho.
Esse diagnóstico abriu uma linha de pesquisa diretamente relevante para este trabalho,
explorada em~\cite{Dickens:10}, \cite{LiaoIPDPS:07}, \cite{Gropp:08} e~\cite{Liao:08},
e fundamenta a escolha de \texttt{MPI\_File\_write\_at} como primitiva de E/S paralela
no SDDP.

\citeauthor{LiaoIPDPS:07}~\cite{LiaoIPDPS:07} demonstraram que escritas
independentes cujo tamanho não preenche uma unidade de \textit{stripe} do
Lustre forçam a aquisição de múltiplos locks sequenciais, degradando o
desempenho. Como solução, os autores propuseram um mecanismo de
\textit{two-stage write-behind buffering} que acumula dados em memória até
completar um bloco alinhado à fronteira de \textit{stripe} antes de emitir a
escrita ao filesystem. Esta dissertação \emph{não adotou} essa abordagem: a implementação descrita no Capítulo~\ref{cap:proposta}
apenas agrega os registros horários por par
$(\text{estágio}, \text{cenário})$ em uma única escrita, sem
alinhamento explícito às fronteiras de \textit{stripe}. Ainda assim,
parte das escritas resultantes \emph{acaba se alinhando
incidentalmente} ao tamanho de \textit{stripe} configurado no Lustre
(1~MB nos experimentos), em decorrência do padrão de acesso do SDDP
no formato horário: arquivos com algumas centenas de elementos do
sistema elétrico produzem blocos por par
$(\text{estágio}, \text{cenário})$ próximos ou múltiplos da fronteira
de \textit{stripe} (cf.~Seção~\ref{sec:formato_resultados}). Esse
alinhamento incidental tende a beneficiar parte das operações de E/S
do SDDP --- evitando, para esses blocos, a aquisição de múltiplos
locks descrita por \citeauthor{LiaoIPDPS:07} --- e ajuda a explicar o
desempenho já obtido pela arquitetura proposta mesmo na ausência de
um mecanismo explícito de bufferização alinhada. A incorporação
\emph{sistemática} desse alinhamento permanece como direção de
aprimoramento da arquitetura proposta, beneficiando inclusive os arquivos
cujos blocos por cenário ficam desalinhados.

Em continuação direta, \citeauthor{Dickens:10}~\cite{Dickens:10} propuseram a Y-lib,
uma biblioteca que substitui o driver ADIO padrão do ROMIO para Lustre, implementando
hierarchical striping e split writing para alinhar as operações de E/S às fronteiras
de stripe e reduzir a contenção de locks. Os experimentos reportados pelos autores
mostram ganhos de até 220\% sobre a implementação padrão, confirmando que o Lustre
requer tratamento específico para extrair desempenho em escritas paralelas.

\citeauthor{Gropp:08}~\cite{Gropp:08} formalizaram que, quando os acessos de
diferentes processos não são interleaved --- isto é, cada processo escreve em
uma região exclusiva do arquivo ---, a chamada \texttt{MPI\_File\_write\_all}
degenera para \texttt{MPI\_File\_write} acrescida de overhead de sincronização,
tornando as escritas coletivas estritamente piores do que as independentes nesse
regime. Esse resultado fundamenta teoricamente a escolha de
\texttt{MPI\_File\_write\_at} no SDDP, cujo padrão de acesso é precisamente
não-interleaved.

\citeauthor{Liao:08}~\cite{Liao:08} mostraram que o particionamento de
\textit{file domains} do two-phase I/O do ROMIO não coincide com as fronteiras
de \textit{stripe} do Lustre, gerando contenção de locks mesmo em operações
coletivas bem-formadas. Os autores propuseram um particionamento adaptativo
ao protocolo de locking do filesystem subjacente, reduzindo essa contenção.
Esse resultado complementa \cite{Dickens:08}, reforçando que, para o padrão
de acesso do SDDP, escritas independentes com offsets pré-calculados evitam
inteiramente o problema de particionamento identificado pelos autores.

\citeauthor{GarciaCarballeiraISPDC:23}~\cite{GarciaCarballeiraISPDC:23}
apresentam o \emph{Expand Ad-Hoc}, um sistema de arquivos paralelo
\textit{ad-hoc} baseado em servidores MPI em espaço de usuário que
virtualiza o armazenamento local dos nós de computação --- SSDs ou
memória compartilhada --- em um volume efêmero, cujo ciclo de vida
coincide com o do \textit{job} submetido: a aplicação escreve nesse
volume durante a execução e a sincronização com o \textit{backend}
persistente (Lustre, GPFS) ocorre apenas nas etapas de
\textit{stage-in} e \textit{stage-out}, fora da janela crítica da
computação. Essa abordagem é diretamente relevante para a
degradação observada no Experimento SDumont
(Subseções~\ref{sec:sdumont_atual} e~\ref{sec:sdumont_mpiio}), em
que a concorrência entre o \textit{job} do SDDP e outras cargas de
HPC sobre o mesmo Lustre compartilhado impôs um teto absoluto de
banda (34{,}1~Gb/s em 32~nós) uma ordem de grandeza inferior ao
obtido na AWS com FSx for Lustre dedicado (380{,}3~Gb/s). Compor o
MPI-IO com escritas independentes sobre uma camada como o Expand
Ad-Hoc, instanciada nos próprios nós alocados ao \textit{job} SDDP,
contornaria inteiramente essa contenção externa, com migração ao
Lustre apenas no \textit{stage-out} final, e tem potencial de
aproximar o desempenho observado no SDumont do desempenho obtido em
um \textit{backend} dedicado.